Everything so far warned us. The model knows \texttt{abs} and
\texttt{add} individually. Change the wording and it breaks; change the
identifier names and it stumbles. But the pieces are there. So the
question this chapter asks is the sharpest one of the whole book:

\begin{quote}
If the model knows how to \emph{add} and how to take an \emph{absolute
value}, can it do the thing the two of them make together --- compute
the absolute value of a sum?
\end{quote}

This is \textbf{compositional generalization}: taking knowledge the
model already has and applying it to a structure it never saw during
training. It is arguably the deepest question in machine learning,
because it separates ``the model memorized something'' from ``the model
\emph{understands} something.'' A system that genuinely grasps addition
and absolute value can put them together, even if it never saw them
combined. A system that merely memorized the word ``sum'' mapped to a
\texttt{+} cannot.

\section{The setup: design the test to be
unfair-but-fair}\label{the-setup-design-the-test-to-be-unfair-but-fair}

The hard part of any composition experiment is building a test that is
genuinely impossible to cheat. Here is how we do it.

We train the model exactly as before, on \textbf{primitive} operations
only: \texttt{add}, \texttt{sub}, \texttt{mul}, \texttt{abs}, and the
rest. We never show it a composite operation.

Then we build a set of \textbf{composite} operations that are \emph{made
from} the primitives: \texttt{abs\_add} (absolute value of a sum),
\texttt{sub\_abs} (absolute difference), \texttt{sum\_of\_abs} (sum of
absolute values), \texttt{max\_add}, \texttt{double\_first},
\texttt{len\_minus\_one}. These are not new primitive functions; they
are combinations of the ones the model already knows. To solve
\texttt{abs\_add}, a model \emph{must} perform an addition and then take
an absolute value --- two pieces it demonstrably knows.

The key design decision is that these composites were \textbf{never
shown during training} (the teacher holds them out, exactly as it holds
out test tasks). So the test is: can the model reach into its bag of
known primitive behaviors and compose them in a combination it has never
seen? If it cannot, we have found its compositional boundary. If it can,
we have found something remarkable.

To make the test fair, we also show what \emph{both} baselines do. The
retrieval baseline has nothing to compose --- it can only find the
closest primitive example and copy it. The template baseline has no rule
for \texttt{abs\_add} either, because the teacher never gave it one. So
if \emph{everything} fails, we have a fair statement: composition cannot
be achieved by the easy shortcuts, and the model too cannot do it.

\section{What we measured}\label{what-we-measured}

Here is the result, and it is unambiguous:

\begin{longtable}[]{@{}
  >{\raggedright\arraybackslash}p{0.1624\linewidth}
>{\raggedright\arraybackslash}p{0.1894\linewidth}
>{\raggedright\arraybackslash}p{0.2706\linewidth}
>{\raggedright\arraybackslash}p{0.2976\linewidth}
@{}}
\toprule\noalign{}
\begin{minipage}[b]{\linewidth}\raggedright
Test
\end{minipage} & \begin{minipage}[b]{\linewidth}\raggedright
MODEL
\end{minipage} & \begin{minipage}[b]{\linewidth}\raggedright
TEMPLATE
\end{minipage} & \begin{minipage}[b]{\linewidth}\raggedright
RETRIEVAL
\end{minipage} \\
\midrule\noalign{}
\endhead
\bottomrule\noalign{}
\endlastfoot
In-distribution (primitives) & 0.900 & 0.800 & 1.000 \\
\textbf{Compositional (unseen combos)} & \textbf{0.000} & \textbf{0.000}
& \textbf{0.000} \\
\end{longtable}

Let us read this carefully, because the meaning is subtle.

On the primitive operations it saw, the model scores 0.900 --- good. But
on the composite operations, it scores \textbf{0.000}. Not 0.1, not
``sometimes,'' not ``close.'' Zero. The model never produced a correct
program for any composite operation, even though every composite is a
combination of primitives it individually knows.

And here is the crucial fairness check: the baselines fail too. The
retrieval baseline gets 0.000 because it has nothing to retrieve ---
there is no composite example in its memory. The template baseline gets
0.000 because no template rule covers a composite. So this is not a case
where the model is merely \emph{worse} than an easy shortcut. This is a
case where the easy shortcuts also collapse, which means the only way to
succeed is to \emph{genuinely compose} --- and the model cannot.

\section{What this really means}\label{what-this-really-means}

This is the cleanest negative result in the book, and it is the book's
central finding. Let us unpack it.

The model knows \texttt{abs} at 0.900 and \texttt{add} at 0.900. But
\texttt{add} the absolute values: 0.000. It is as if the model has two
perfectly good rooms --- one with an ``add'' skill, one with an ``abs''
skill --- but no door between them. It cannot take the output of one and
feed it into the other, even though that is the simplest possible
combination.

This is the deep meaning of ``surface-tied'' that we began to glimpse in
Chapter 15. It is not merely that the model is sensitive to wording. It
is that the model's knowledge is \textbf{flat and associative}, not
\textbf{structured and compositional}. It learned that the token
sequence for ``sum'' predicts the token sequence for \texttt{+}. It did
not learn a general \emph{operation} that can be an argument to another
general operation. The knowledge is a set of reflexes, not a toolbox
with reusable parts.

This distinction is why we built the four-level generalization taxonomy
in the first place and why we refused to collapse everything into one
number. If we had only reported ``the model scores 0.900 on this task,''
we would have given a genuinely misleading picture. The model is
excellent at what it was shown --- and utterly unable to combine what it
was shown into anything new. Both statements are true, and both must be
reported.

\section{The puzzling asymmetry (and why it's a finding, not a
bug)}\label{the-puzzling-asymmetry-and-why-its-a-finding-not-a-bug}

One thing will strike us: the model scores 0.900 in-distribution but
0.000 on composition. Why would a model that is ``good'' fail so
completely on something so similar?

The answer is that the two numbers measure \emph{different} kinds of
capability, which is precisely the point of the taxonomy. The 0.900
measures \textbf{distributional generalization} --- new examples of the
\emph{same} operation. The 0.000 measures \textbf{structural
generalization} --- new \emph{constructions} combining known operations.
These are different abilities, and a model can be high on one and zero
on the other. That asymmetry is not a contradiction; it is the shape of
the truth we are mapping.

\section{Why we don't call it a ``failure of one thing'' without
checking the other
two}\label{why-we-dont-call-it-a-failure-of-one-thing-without-checking-the-other-two}

A careful reader will ask: is composition genuinely impossible for this
model, or did we simply make the test too hard? The baselines answer
this. If the test were simply ``too hard'' in some absolute sense, the
retrieval and template baselines would not necessarily fail --- but they
are the \emph{easy} methods, and they fail too. So the test is not too
hard for an easy method; it is unsolvable by any method except genuine
composition, and the model has no composition. This is a fair,
controlled comparison, which is exactly what our one-variable discipline
demands.

\section{The shortcut audit}\label{the-shortcut-audit}

Before we trust ``0.000,'' we ask how the model might have \emph{faked}
a success.

\begin{enumerate}
\def\labelenumi{\arabic{enumi}.}
\tightlist
\item
  \textbf{Suspected shortcut: memorizing a composite example.} Ruled out
  --- no composite operation was ever shown during training, so there is
  nothing to memorize.
\item
  \textbf{Suspected shortcut: retrieving a near-identical primitive and
  patching it.} Ruled out --- the retrieval baseline, which does exactly
  this, scores 0.000.
\item
  \textbf{Suspected shortcut: a rule that ``just works'' on composites.}
  Ruled out --- the template baseline, which has rule-based generation,
  scores 0.000.
\item
  \textbf{Suspected shortcut: the judge being too strict or too loose.}
  This is the one we genuinely have to hold onto as a question, and it
  is why the judge is mutation-tested (Chapter 17). If the judge were
  rejecting correct composite programs, we might be misreading a real
  success as a failure.
\end{enumerate}

\section{The honest framing}\label{the-honest-framing}

This is a \emph{useful}, \emph{clean} negative result --- and in
science, a decisive negative is often more informative than a vague
positive. We now know something sharp about tiny models: \textbf{at
\textasciitilde1M parameters, they do not compose.} They learn reflexes,
not structure. This is not a statement about ``intelligence'' or about
``what models will eventually do.'' It is a precise, reproducible
measurement of where this model's capability stops.

And it sets up the rest of the book in a way that is far more
interesting than if the model had composed successfully. Because now the
question becomes: \emph{what would it take to give a tiny model this
ability?} More capacity? More data? Different data with more structure?
That is what the next chapters investigate --- and the answer is one of
the most genuinely surprising things the book finds.
